\SourceLeaf{026}
\Article{14}
\textsc{Lemma.} \textit{Sind die Größe $r$ und der Winkel $\varphi$ so bestimmt, dass die Gleichungen}
\[
\begin{aligned}
&r^m\cos m\varphi+Ar^{m-1}\cos(m-1)\varphi+Br^{m-2}\cos(m-2)\varphi+\text{usw.}\\
&\hspace{3em}+Krr\cos2\varphi+Lr\cos\varphi+M=0\quad.\quad.\quad[1]\\
&r^m\sin m\varphi+Ar^{m-1}\sin(m-1)\varphi+Br^{m-2}\sin(m-2)\varphi+\text{usw.}\\
&\hspace{3em}+Krr\sin2\varphi+Lr\sin\varphi=0\quad.\quad.\quad[2]
\end{aligned}
\]
\textit{gelten, so ist das Polynom $x^m+Ax^{m-1}+Bx^{m-2}+\text{usw.}+Kxx+Lx+M=X$ durch den quadratischen Faktor $xx-2\cos\varphi\mathbin{.}rx+rr$ teilbar, sofern nicht $r\sin\varphi$ $=0$ gilt; ist dagegen $r\sin\varphi=0$, so ist dasselbe Polynom durch den linearen Faktor $x-r\cos\varphi$ teilbar.}

\textit{Beweis.} I. Nach dem vorigen Artikel sind alle folgenden Größen durch $xx-2\cos\varphi\mathbin{.}rx+rr$ teilbar:
{\small
\[
\setlength{\arraycolsep}{2pt}
\begin{array}{rll}
\sin\varphi\mathbin{.}rx^m&-\sin m\varphi\mathbin{.}r^m x&+\sin(m-1)\varphi\mathbin{.}r^{m+1}\\
A\sin\varphi\mathbin{.}rx^{m-1}&-A\sin(m-1)\varphi\mathbin{.}r^{m-1}x&+A\sin(m-2)\varphi\mathbin{.}r^m\\
B\sin\varphi\mathbin{.}rx^{m-2}&-B\sin(m-2)\varphi\mathbin{.}r^{m-2}x&+B\sin(m-3)\varphi\mathbin{.}r^{m-1}\\
&\text{usw.}&\text{usw.}\\
K\sin\varphi\mathbin{.}rxx&-K\sin2\varphi\mathbin{.}rrx&+K\sin\varphi\mathbin{.}r^3\\
L\sin\varphi\mathbin{.}rx&-L\sin\varphi\mathbin{.}rx&\qquad *\\
M\sin\varphi\mathbin{.}r&\qquad *&+M\sin(-\varphi)\mathbin{.}r
\end{array}
\]
}
Daher ist auch ihre Summe durch $xx-2\cos\varphi\mathbin{.}rx+rr$ teilbar. Die ersten Teile dieser Größen bilden aber die Summe $\sin\varphi\mathbin{.}rX$; die zweiten ergeben zusammen nach $[2]$ den Wert $0$. Dass auch die Summe der dritten Teile verschwindet, erkennt man leicht, wenn man $[1]$ mit $\sin\varphi$ und $[2]$ mit $\cos\varphi$ multipliziert und das erste Produkt vom zweiten abzieht. Daraus folgt, dass das Polynom $\sin\varphi\mathbin{.}rX$ durch $xx-2\cos\varphi\mathbin{.}rx+rr$ teilbar ist, und damit auch das Polynom $X$, sofern nicht $r\sin\varphi=0$ gilt. Q. E. P.

II. Ist dagegen $r\sin\varphi=0$, so ist entweder $r=0$ oder $\sin\varphi=0$. Im ersten Fall ist nach $[1]$ auch $M=0$, und daher ist $X$ durch $x$ oder durch $x-r\cos\varphi$ teilbar; im zweiten Fall ist $\cos\varphi=\pm1$, $\cos2\varphi=+1$, $\cos3\varphi=\pm1$ und allgemein $\cos n\varphi=\cos\varphi^n$. Wegen $[1]$ ergibt sich also $X=0$, wenn man $x=r\cos\varphi$ setzt, und folglich ist das Polynom $X$ durch $x-r\cos\varphi$ teilbar. Q. E. S.

\SourceLeaf{027}
\Article{15}
Der vorige Satz wird gewöhnlich mit Hilfe imaginärer Größen bewiesen; siehe \textsc{Euler}, \textit{Introd. in Anal. Inf.}, Bd. I, S. 110. Ich hielt es für der Mühe wert zu zeigen, wie er sich ebenso leicht ohne ihre Hilfe herleiten lässt. Es ist nun klar, dass zum Beweis unseres Satzes nur noch Folgendes zu zeigen ist: \textit{Für jedes vorgegebene Polynom $X$ der Form $x^m+Ax^{m-1}+Bx^{m-2}+\text{usw.}+Lx+M$ lassen sich $r$ und $\varphi$ so bestimmen, dass die Gleichungen $[1]$ und $[2]$ gelten.} Daraus folgt nämlich, dass $X$ einen reellen Faktor ersten oder zweiten Grades besitzt. Die Division ergibt notwendig einen reellen Quotienten niedrigeren Grades, der aus demselben Grund ebenfalls einen Faktor ersten oder zweiten Grades besitzt. Setzt man dieses Verfahren fort, so wird $X$ schließlich in reelle lineare oder quadratische Faktoren zerlegt. Diesen Satz zu beweisen ist also das Ziel der folgenden Untersuchungen.

\Needspace{15\baselineskip}
\Article{16}
Man denke sich eine feste unendliche Ebene (die Ebene der Tafel, Fig. 1) und darin eine feste unendliche Gerade $GC$, die durch einen festen Punkt $C$ geht. Man wähle eine Länge als Einheit, sodass sich alle Geradenlängen durch Zahlen ausdrücken lassen. In jedem Punkt $P$ der Ebene, dessen Abstand vom Mittelpunkt $C$ gleich $r$ ist und dessen Winkel $GCP=\varphi$ beträgt, errichte man eine Senkrechte gleich dem Wert des Ausdrucks
\[
r^m\sin m\varphi+Ar^{m-1}\sin(m-1)\varphi+\text{usw.}+Lr\sin\varphi
\]
den ich im Folgenden der Kürze halber stets mit $T$ bezeichnen werde. Den Abstand $r$ betrachte ich stets als positiv. Für Punkte auf der anderen Seite der Achse muss der Winkel $\varphi$ entweder als größer als zwei rechte Winkel oder als negativ angesehen werden (was hier auf dasselbe hinausläuft). Die Endpunkte dieser Senkrechten (die für einen positiven Wert von $T$ oberhalb der Ebene, für einen negativen unterhalb und bei verschwindendem Wert in der Ebene selbst zu nehmen sind) liegen auf einer stetigen gekrümmten Fläche, die sich nach allen Richtungen ins Unendliche erstreckt. Ich werde sie im Folgenden der Kürze halber die \textit{erste Fläche} nennen. Ganz ebenso beziehe man auf dieselbe Ebene, denselben Mittelpunkt und dieselbe Achse eine andere Fläche, deren Höhe über jedem Punkt der Ebene gleich
\[
r^m\cos m\varphi+Ar^{m-1}\cos(m-1)\varphi+\text{usw.}+Lr\cos\varphi+M
\]
ist; diesen Ausdruck werde ich der Kürze halber stets mit $U$ bezeichnen. Diese Fläche, die ebenfalls stetig ist und sich nach allen Richtungen ins Unendliche erstreckt, werde ich von der vorigen durch die Bezeichnung
\SourceLeaf{028}
\textit{zweite Fläche} unterscheiden. Damit ist klar, dass die ganze Aufgabe darin besteht zu zeigen, dass es mindestens einen Punkt gibt, der zugleich in der Ebene, auf der ersten Fläche und auf der zweiten Fläche liegt.

\Article{17}
Man erkennt leicht, dass die erste Fläche teils oberhalb, teils unterhalb der Ebene liegt. Denn offenbar lässt sich der Abstand $r$ vom Mittelpunkt so groß wählen, dass die übrigen Glieder in $T$ gegenüber dem ersten Glied $r^m\sin m\varphi$ verschwinden; dieses kann aber durch geeignete Wahl des Winkels $\varphi$ sowohl positiv als auch negativ werden. Die feste Ebene wird daher notwendig von der ersten Fläche geschnitten. Diese Schnittlinie der Ebene mit der ersten Fläche werde ich die \textit{erste Linie} nennen; sie wird also durch die Gleichung $T=0$ bestimmt. Aus demselben Grund wird die Ebene von der zweiten Fläche geschnitten; der Schnitt bildet eine durch die Gleichung $U=0$ bestimmte Kurve, die ich die \textit{zweite Linie} nennen werde. Eigentlich besteht jede der beiden Kurven aus mehreren Zweigen, die vollständig voneinander getrennt sein können, während jeder für sich eine stetige Linie ist. Die erste Linie ist sogar stets eine sogenannte zusammengesetzte Kurve, und die Achse $GC$ ist als Teil dieser Kurve anzusehen: Welchen Wert man $r$ auch gibt, es gilt stets $T$ $=0$, wenn $\varphi$ $=0$ oder $=180^{\circ}$ ist. Doch ist es besser, die Gesamtheit aller Zweige, die durch sämtliche Punkte mit $T=0$ gehen, als eine einzige Kurve zu betrachten (entsprechend dem in der höheren Geometrie allgemein üblichen Gebrauch), und ebenso sämtliche Zweige durch alle Punkte mit $U=0$. Es ist nun klar, dass die Aufgabe darauf zurückgeführt ist zu zeigen, dass es mindestens einen Punkt in der Ebene gibt, in dem irgendein Zweig der ersten Linie von einem Zweig der zweiten Linie geschnitten wird. Dazu müssen wir die Beschaffenheit dieser Linien näher betrachten.

\Article{18}
Zunächst bemerke ich, dass beide Kurven algebraisch sind, und zwar bei Darstellung durch rechtwinklige Koordinaten vom Grad $m$. Wählt man nämlich den Anfangspunkt der Abszissen in $C$, die Abszissen $x$ in Richtung $G$ und die Ordinaten $y$ in Richtung $P$, so gilt $x=r\cos\varphi$, $y=r\sin\varphi$ und daher allgemein, welchen Wert $n$ auch haben mag,
\[
\begin{aligned}
r^n\sin n\varphi={}&nx^{n-1}y-\frac{n\mathbin{.}n-1\mathbin{.}n-2}{1\mathbin{.}2\mathbin{.}3}x^{n-3}y^3\\
&+\frac{n\mathbin{.}\mathbin{.}\mathbin{.}n-4}{1\mathbin{.}\mathbin{.}\mathbin{.}\mathbin{.}5}x^{n-5}y^5-\text{usw.},\\
r^n\cos n\varphi={}&x^n-\frac{n\mathbin{.}n-1}{1\mathbin{.}2}x^{n-2}yy\\
&+\frac{n\mathbin{.}n-1\mathbin{.}n-2\mathbin{.}n-3}{1\mathbin{.}2\mathbin{.}3\mathbin{.}4}x^{n-4}y^4-\text{usw.}
\end{aligned}
\]
Also bestehen sowohl $T$ als auch $U$ aus mehreren Gliedern der Form $ax^{\alpha}y^{\beta}$,
\SourceLeaf{029}
wobei $\alpha$ und $\beta$ positive ganze Zahlen bezeichnen, deren Summe im größten Fall $=m$ ist. Ferner lässt sich leicht voraussehen, dass alle Glieder von $T$ den Faktor $y$ enthalten und die erste Linie somit eigentlich aus einer Geraden (mit der Gleichung $y=0$) und einer Kurve vom Grad $m-1$ zusammengesetzt ist; doch brauchen wir auf diese Unterscheidung hier nicht einzugehen.

Wichtiger wird die Untersuchung sein, ob die erste und die zweite Linie ins Unendliche verlaufende Arme besitzen und wie viele und welcher Art diese sind. In unendlicher Entfernung vom Punkt $C$ fällt die erste Linie, deren Gleichung $\sin m\varphi+\frac{A}{r}\sin(m-1)\varphi+\frac{B}{rr}\sin(m-2)\varphi\ \text{usw.}=0$ lautet, mit der Linie zusammen, deren Gleichung $\sin m\varphi=0$ ist. Diese stellt $m$ Geraden dar, die sich im Punkt $C$ schneiden. Die erste ist die Achse $GCG'$; die übrigen sind gegen sie unter den Winkeln $\frac{1}{m}180$, $\frac{2}{m}180$, $\frac{3}{m}180$ usw. Grad geneigt. Die erste Linie besitzt daher $2m$ unendliche Zweige, welche den Umfang eines mit unendlichem Radius beschriebenen Kreises in $2m$ gleiche Teile zerlegen: Der erste Zweig schneidet den Umfang im Schnittpunkt von Kreis und Achse, der zweite im Abstand $\frac{1}{m}180^{\circ}$, der dritte im Abstand $\frac{2}{m}180^{\circ}$ usw. Ebenso besitzt die zweite Linie in unendlicher Entfernung vom Mittelpunkt eine durch die Gleichung $\cos m\varphi=0$ ausgedrückte Asymptote. Diese ist eine Gesamtheit von $m$ Geraden, die sich ebenfalls im Punkt $C$ unter gleichen Winkeln schneiden, jedoch so, dass die erste mit der Achse $CG$ den Winkel $\frac{1}{m}90^{\circ}$ bildet, die zweite den Winkel $\frac{3}{m}90^{\circ}$, die dritte den Winkel $\frac{5}{m}90^{\circ}$ usw. Daher besitzt auch die zweite Linie $2m$ unendliche Zweige, von denen jeder die Mitte zwischen zwei benachbarten Zweigen der ersten Linie einnimmt. Sie schneiden also den Umfang eines mit unendlich großem Radius beschriebenen Kreises in Punkten, deren Winkelabstände von der Achse $\frac{1}{m}90^{\circ}$, $\frac{3}{m}90^{\circ}$, $\frac{5}{m}90^{\circ}$ usw. betragen. Ferner ist klar, dass die Achse selbst stets zwei unendliche Zweige der ersten Linie bildet, nämlich den ersten und den an der Stelle $m+1$. Sehr deutlich zeigt sich diese Lage der Zweige in Fig. 2, die für den Fall $m=4$ gezeichnet ist. Die Zweige der zweiten Linie sind dort zur Unterscheidung von denen der ersten punktiert dargestellt; dasselbe gilt auch für die vierte Figur\HistoricalNote{8}. --- Da diese Folgerungen aber von größter Bedeutung sind und unendlich große Größen bei manchen Lesern Anstoß erregen könnten, werde ich im nächsten Artikel zeigen, wie man sie auch ohne die Hilfe von Unendlichkeiten herleiten kann.
